\Needspace{6\baselineskip}
\SourceLeaf{341}

\begin{center}
[3.]
\end{center}

Finally we wish to show how to pass from the proof in [1] to that in [2], since
the proof was in fact first found by means of laterals and only later by geometrical
considerations.

Let us see what form the proof in [1] takes if one point of the quadrilateral, for\,example (4), coincides with
the origin from which we reckon. We then have $p^{\mathrm{IV}} = 0$, so the
coordinates are
\[
\begin{array}{rl}
\text{of I:} & p'p'' \\
\text{of II:} & p'p''' \\
\text{of III:} & p''p'''.
\end{array}
\]
Hence by [2] the sides are
\[
\begin{array}{rl}
(\text{I}, \text{II}) & = (4,1) . (4,2) \\
(\text{II}, \text{III}) & = (4,1) . (4,3) \\
(\text{III}, \text{I}) & = (4,2) . (4,3).
\end{array}
\]
Consequently the sides of the triangle are in the ratio
\[
(\text{I}, \text{II}) : (\text{II}, \text{III}) : (\text{III}, \text{I}) = (4,1) . (4,2) :
(4,1) . (4,3) : (4,2) . (4,3),
\]
as we determined them by construction in [2].
